\section{Measurement audit and reporting checklist}
\label{sec:pitfalls}
A lower classifier rate can change because the model, the reference, or the measurement path changed. Our audit separates these possibilities rather than adding them into a single corrected safety score. The checklist records each trigger, observed consequence, detection method, and practical control; the detailed historical comparisons remain in Appendix~\ref{app:pitfalls}.

\paragraph{The judge's construct and input matter.} HarmBench conditions on a behavior text, so an unverified nearest-behavior mapping changes the question it answers. The historical WildJailbreak top-1 and top-2 references changed labels on 36.65\% of baseline rows. This is reference sensitivity, not an estimate of mapping error. The redesigned SORRY evaluation supplies the original English instruction instead. Even then, fulfillment, reply harmfulness, and nonrefusal remain different constructs. The human-reference confusion matrices expose class- and style-specific errors; an overall agreement score alone cannot show that errors cancel in learned-minus-random comparisons.

\paragraph{Length caps and decoding paths define the comparison.} A short refusal-like prefix can later turn into compliance, and a warning can precede substantive content. Changing the length cap or scoring only the first sentence can therefore change the outcome without changing a stored prefix. We preserve the 2400-token harmful-generation cap, retain the 120-token public-cost protocol as a limitation, and report the registered shorter-reply sensitivity separately. Prefix families give v2 only the continuation while model judges receive the complete prefilled reply. Batch composition, padding, and kernel choices can also change text. The engineering label screen is not proof of equivalence; within-cell paths stay fixed and the single-row contrast regression is separately registered.

\paragraph{Training variation and reference choice need different controls.} Repeated old-recipe seed3 training produced near-identical vectors while changing a few val70 labels. A high cosine similarity is an implementation check, not a guarantee of item-label stability. Random directions answer a different question: how much of a judge-rate change occurs under a perturbation of the same norm without that training recipe? We preserve both controls and do not interpret one seed's rank as a population property.

\paragraph{External validation has a domain.} A released human test split provides a reference for judge implementation and class-specific errors, not human labels for our protected-model outputs. A split can contain different replies to related instructions, and translation-style replies can already be English. We report those properties alongside accuracy and $\kappa$. No judge is discarded after an unfavorable protected-model effect, and low-positive-count encoding styles are not called validated merely because accuracy is high.

\paragraph{An audit trail must distinguish decisions from receipts.} Source and model hashes, row-level atomic records, runtime vector norms, and cache-source hashes make implementation changes traceable. They do not repair an incorrect construct. The first nonzero lease preceded its dedicated release receipt; the prerequisites were already frozen, but the receipt is explicitly retrospective. Search early-stop scores likewise remain internal search signals rather than independent fulfillment judgments. Failed criteria, aborted work, and incomplete model ports stay in the record instead of being rewritten as successful preregistered tests.
